\documentclass{article}

\usepackage[margin=1in]{geometry}

\title{Mock FF7 Style Project \\ Snapshot Objectives}
\author{Team 13}
\date{Spring 2026}

\begin{document}

\maketitle

\tableofcontents
\newpage

\section{Introduction}

This document explains the goals and progress for each project snapshot
throughout the development of the Mock FF7 Style Project.
Each snapshot represents a checkpoint where the team adds new features,
updates documentation, and improves the overall structure of the project.

\section{Snapshot 1 Starting the Project}

\subsection{Main Goals}

For the first snapshot, the main goal was to set up the project and create
the basic structure needed for development.

The team focused on:

\begin{itemize}
\item Creating the GitHub repository
\item Setting up the Jira board for sprint planning
\item Building the first version of the combat system
\item Creating the initial SDD and SRS documents
\item Setting up Docker files
\item Organizing the project structure
\item Creating the user manual
\end{itemize}

\subsection{Expected Outcome}

By the end of Snapshot 1, the project should have a working
combat prototype along with the required documentation and setup files.

\section{Snapshot 2 First Checkpoint}

\subsection{Reflection on Snapshot 1}

During Snapshot 1, the team successfully created the repository,
set up the development environment, and built the first version
of the RPG combat system.

\subsection{New Goals}

For Snapshot 2, the team plans to improve the gameplay systems
and continue expanding the documentation.

The new goals include:

\begin{itemize}
\item Adding an inventory system
\item Improving combat balance
\item Adding potion and healing features
\item Creating TestRail testing documents
\item Updating the workflow diagram
\item Expanding the SDD and SRS documents
\item Adding more Jira sprint tasks
\end{itemize}

\subsection{Expected Outcome}

At the end of Snapshot 2, the project should feel more complete
and include additional gameplay features along with testing documentation.

\section{Snapshot 3 Second Checkpoint}

\subsection{Reflection on Snapshot 2}

Snapshot 2 added more gameplay systems and improved the structure
of the project. The team also completed testing documentation
and updated the workflow design.

\subsection{New Goals}

The goals for Snapshot 3 are focused on improving the quality
and organization of the project.

These goals include:

\begin{itemize}
\item Expanding combat actions and abilities
\item Improving code organization
\item Adding more testing coverage
\item Improving Docker configuration
\item Updating all project documentation
\item Continuing sprint planning in Jira
\end{itemize}

\subsection{Expected Outcome}

By the end of Snapshot 3, the project should be more stable,
better organized, and closer to the final version.

\section{Snapshot 4 Final Submission}

\subsection{Reflection on Snapshot 3}

Snapshot 3 improved the overall structure of the project
and expanded the gameplay systems while continuing testing
and documentation updates.

\subsection{Final Goals}

The final snapshot focuses on polishing the project
and completing all remaining work.

The final goals include:

\begin{itemize}
\item Finalizing all project documents
\item Completing TestRail reports
\item Cleaning up the repository structure
\item Finalizing Docker setup
\item Updating workflow diagrams
\item Adding future improvement ideas
\item Completing the final project review
\end{itemize}

\subsection{Future Improvements}

Some features the team would like to add in the future include:

\begin{itemize}
\item Multiplayer support
\item Save and load functionality
\item A graphical user interface
\item More enemy types
\item Character leveling and progression
\end{itemize}

\subsection{Expected Outcome}

By the end of Snapshot 4, the project should include
all required documentation, testing reports, and a finalized
mock RPG prototype ready for submission.

\section{Conclusion}

These snapshots help track the progress of the Mock FF7 Style Project
throughout development. Each checkpoint allows the team to improve
the project step by step while practicing software engineering workflows,
documentation, testing, and deployment planning.

\end{document}
```
